% ============================================================
% PROPOSAL PROYEK AKHIR — FARCHAN DEANO MUHAMMAD (CAN)
% BAB 2 — KAJIAN PUSTAKA (RINGKAS)
%
% Project: Sketchbook Universe
% Scope  : Frontend / HITL UI / Maskot Momo / Simulasi Interaktif Berlevel
% Format : PENS Standard (A4, 12pt TNR, 1.5 spacing, 4cm/3cm margin)
% Style  : Mengikuti pola format Izzah (PilAItes 2025) — RINGKAS, ngalir
%
% File ini adalah STANDALONE COMPILABLE LaTeX file berisi Bab 2 saja.
% Target halaman: 8--10 halaman (RINGKAS, detail di Bab 3).
%
% Struktur:
%   2.1 Deskripsi Permasalahan          (~1 hal)
%   2.2 Teori Penunjang                  (~5--6 hal, 7 sub-bab)
%     2.2.1 Literasi Kecerdasan Buatan
%     2.2.2 Human-in-the-Loop
%     2.2.3 Explainable Interface dan Confidence Score
%     2.2.4 Simulasi Interaktif Berlevel [GBL]
%     2.2.5 Pedagogical Agent --- Maskot Momo
%     2.2.6 Child-Computer Interaction untuk SMP
%     2.2.7 Usability Testing
%   2.3 Penelitian Terkait               (~2--3 hal, 5 sub-bab + SOTA)
%     2.3.1 AI Literacy in K-12
%     2.3.2 HITL Machine Learning
%     2.3.3 Game-Based Learning in CS Education
%     2.3.4 Children Behavioral Patterns in AI Learning
%         [POST-PIVOT: gantikan Real-Time Hand Gesture yang DEPRECATED]
%     2.3.5 State of The Art (tabel POLOSAN, 6 sistem)
%
% Catatan post-pivot 16/6/26:
%   - Sub-bab MediaPipe (B02 Meng) DIHAPUS dari teori penunjang.
%   - Sub-bab Real-Time Hand Gesture Monitoring DIHAPUS, diganti
%     Children Behavioral Patterns in AI Learning (B07 Hsu, B08 Ocak).
%   - B02 tetap di Daftar Pustaka dengan catatan DEPRECATED
%     (dipertahankan untuk dokumentasi historis canvas drawing input).
% ============================================================

\documentclass[12pt,a4paper,oneside,openany]{report}

% ============================================================
% PACKAGES
% ============================================================
\usepackage[top=4cm,left=4cm,right=3cm,bottom=3cm]{geometry}
\usepackage{graphicx}
\usepackage{booktabs}
\usepackage{array}
\usepackage{tabularx}
\usepackage{longtable}
\usepackage{float}
\usepackage{caption}
\usepackage{titlesec}
\usepackage{indentfirst}
\usepackage{setspace}
\usepackage{etoolbox}
\usepackage{fancyhdr}
\usepackage{afterpage}
\usepackage{amsmath}
\usepackage{amssymb}
\usepackage{url}
\usepackage{hyperref}
\hypersetup{colorlinks=false, pdfborder={0 0 0}}

% Times New Roman (pdflatex standard)
\usepackage{times}
\renewcommand{\rmdefault}{ptm}

% Line spacing 1.5
\setstretch{1.5}

% Paragraph indent (1.25 cm, PENS standard)
\setlength{\parindent}{1.25cm}
\setlength{\parskip}{0pt}

% Chapter / section / subsection formatting (pola Izzah)
\titleformat{\chapter}[display]
  {\normalfont\fontsize{14pt}{16.8pt}\selectfont\bfseries\centering}
  {BAB \thechapter}
  {12pt}
  {\MakeUppercase}

\titleformat{\section}
  {\normalfont\fontsize{12pt}{14.4pt}\selectfont\bfseries}
  {\thesection}
  {1em}
  {\MakeUppercase}

\titleformat{\subsection}
  {\normalfont\fontsize{12pt}{14.4pt}\selectfont\bfseries}
  {\thesubsection}
  {1em}
  {}

\titlespacing*{\chapter}{0pt}{0pt}{24pt}
\titlespacing*{\section}{0pt}{18pt}{12pt}
\titlespacing*{\subsection}{0pt}{14pt}{8pt}

% Caption formatting (pola Izzah)
\captionsetup[table]{font=small,labelfont=bf,name={Tabel},labelsep=period,justification=centering,aboveskip=6pt,belowskip=0pt,position=top}
\captionsetup[figure]{font=small,labelfont=bf,name={Gambar},labelsep=period,justification=centering,aboveskip=0pt,belowskip=6pt,position=bottom}

% Header / footer (page number top-right)
\pagestyle{fancy}
\fancyhf{}
\fancyhead[R]{\thepage}
\renewcommand{\headrulewidth}{0pt}
\renewcommand{\footrulewidth}{0pt}

\fancypagestyle{plain}{
  \fancyhf{}
  \fancyhead[R]{\thepage}
  \renewcommand{\headrulewidth}{0pt}
}

% TOC naming
\renewcommand{\contentsname}{DAFTAR ISI}
\renewcommand{\listfigurename}{DAFTAR GAMBAR}
\renewcommand{\listtablename}{DAFTAR TABEL}

% ============================================================
\begin{document}
% ============================================================

% Cover page (minimal, for standalone compile)
\begin{titlepage}
\centering
\vspace*{2cm}
{\fontsize{14pt}{16pt}\selectfont PROPOSAL PROYEK AKHIR}\\[1cm]
{\fontsize{14pt}{18pt}\selectfont\bfseries\MakeUppercase{Pengembangan Simulasi Interaktif Berlevel dengan Mekanisme Human-in-the-Loop pada Sistem Literasi Kecerdasan Buatan untuk Siswa SMP}}\\[2cm]
{\fontsize{12pt}{14pt}\selectfont Oleh:}\\[0.3cm]
{\fontsize{12pt}{14pt}\selectfont\bfseries Farchan Deano Muhammad}\\[0.2cm]
{\fontsize{12pt}{14pt}\selectfont NRP. 53226000xx}\\[2cm]
{\fontsize{12pt}{14pt}\selectfont PROGRAM STUDI SARJANA TERAPAN TEKNOLOGI REKAYASA MULTIMEDIA}\\[0.2cm]
{\fontsize{12pt}{14pt}\selectfont JURUSAN TEKNOLOGI MULTIMEDIA DAN JARINGAN}\\[0.2cm]
{\fontsize{12pt}{14pt}\selectfont POLITEKNIK ELEKTRONIKA NEGERI SURABAYA}\\[0.2cm]
{\fontsize{12pt}{14pt}\selectfont 2025}
\end{titlepage}

% ============================================================
% BAB 2 — KAJIAN PUSTAKA (RINGKAS)
% Citation numbering: skema Sitasi_Final_v1.md (Can, 32 refs)
% ============================================================
%% BEGIN BAB 2 CONTENT %%
\chapter{KAJIAN PUSTAKA}
\label{chap:bab2}

Bab ini menyajikan kajian pustaka yang menjadi landasan teoretis bagi perancangan sistem \textit{Sketchbook Universe}. Kajian dimulai dari deskripsi permasalahan yang menjelaskan mengapa media yang ada saat ini belum memadai untuk memfasilitasi interaksi reflektif siswa SMP terhadap keluaran kecerdasan buatan, dilanjutkan dengan teori penunjang yang mendasari setiap keputusan desain, serta ditutup dengan penelitian terkait dan analisis \textit{state of the art} yang memosisikan proyek ini di antara kajian-kajian sebelumnya. Pembahasan teori pada bab ini bersifat ringkas; rincian implementasi teknis akan diuraikan pada Bab III.

% ============================================================
\section{DESKRIPSI PERMASALAHAN}
\label{sec:deskripsi-permasalahan}
% ============================================================

Perkembangan kecerdasan buatan (KB) telah mengubah cara masyarakat berinteraksi dengan teknologi digital. Siswa SMP kelas 7--9, yang berada dalam rentang usia 13--15 tahun, merupakan generasi yang tumbuh di lingkungan jenuh sistem berbasis KB---mulai dari rekomendasi konten hingga asisten virtual. Namun, keterpaparan terhadap penggunaan sistem KB tidak serta-merta berbanding lurus dengan kemampuan mengevaluasi keluaran yang dihasilkan oleh sistem tersebut. Ng et al. \cite{ref1} menegaskan bahwa literasi KB bukan sekadar kemampuan menggunakan alat berbasis KB, melainkan mencakup pemahaman tentang bagaimana sistem KB menghasilkan keluaran, mengapa keluaran tersebut bisa keliru, dan bagaimana manusia seharusnya berinteraksi dengan keluaran yang bersifat probabilistik.

Media edukasi KB yang tersedia saat ini umumnya memiliki dua kelemahan utama. Pertama, sebagian besar media bersifat informatif-ekspositif---menyajikan penjelasan konseptual tentang KB tanpa menyediakan momen di mana siswa secara langsung menghadapi keluaran KB yang ambigu atau keliru \cite{ref2}. Kedua, media yang bersifat interaktif seringkali memposisikan KB sebagai \textit{oracle} yang selalu benar, sehingga justru memperkuat kecenderungan penerimaan otomatis (\textit{automation bias}) tanpa memberikan ruang bagi siswa untuk mempertanyakan atau mengkoreksi keluaran yang diberikan \cite{ref4}\cite{ref5}. Kondisi ini bertentangan dengan prinsip \textit{Human-in-the-Loop} (HITL) yang menekankan perlunya keterlibatan manusia sebagai evaluator pada titik-titik keputusan kritis \cite{ref4}.

Selain itu, desain antarmuka untuk siswa SMP memerlukan perhatian khusus terhadap aspek perkembangan kognitif dan motivasi. Casal-Otero et al. \cite{ref2} menemukan bahwa strategi literasi KB untuk kalangan K-12 harus bersifat sesuai usia (\textit{age-appropriate}), berbasis pengalaman langsung, dan menghindari abstraksi yang berlebihan. Siswa SMP membutuhkan media yang konkret, bertahap, dan memberikan umpan balik kontekstual---bukan media yang hanya menyajikan teori tentang KB. Dengan demikian, dibutuhkan sebuah sistem interaktif dengan konsekuensi nyata yang memfasilitasi interaksi reflektif siswa SMP terhadap keluaran KB, di mana siswa diberikan kesempatan untuk mengevaluasi, menolak, atau mengkoreksi keluaran tersebut dalam konteks simulasi yang terstruktur.

% ============================================================
\section{TEORI PENUNJANG}
\label{sec:teori-penunjang}
% ============================================================

Bagian ini membahas teori-teori yang menjadi dasar perancangan sistem \textit{Sketchbook Universe}. Setiap sub-bagian menjelaskan teori secara konseptual dan menghubungkannya dengan keputusan desain yang diambil dalam proyek ini. Pembahasan bersifat ringkas; rincian implementasi teknis akan diuraikan pada Bab III.

% ------------------------------------------------------------
\subsection{Literasi Kecerdasan Buatan}
\label{subsec:literasi-kb}
% ------------------------------------------------------------

Literasi kecerdasan buatan didefinisikan oleh Ng et al. \cite{ref1} sebagai seperangkat kompetensi yang memungkinkan individu untuk memahami, menggunakan, mengevaluasi, dan mengkomunikasikan mengenai sistem KB secara kritis. Kerangka konseptual tersebut mencakup empat dimensi utama: (1) mengetahui dan memahami KB, (2) menggunakan dan menerapkan KB, (3) mengevaluasi dan mencipta KB, serta (4) etika KB. Dimensi evaluatif menjadi sangat relevan dalam konteks proyek ini karena mencakup kemampuan menilai apakah keluaran yang dihasilkan sistem KB layak diterima, perlu dipertanyakan, atau harus ditolak. Aspek penting yang sering terabaikan adalah penalaran probabilistik---pemahaman bahwa keluaran sistem KB modern bersifat probabilistik dan bukan jawaban deterministik \cite{ref1}\cite{ref2}.

Dalam konteks K-12, Wang \& Lester \cite{ref11} menegaskan bahwa literasi KB harus diperkenalkan sejak usia sekolah melalui pendekatan yang sesuai dengan tahap perkembangan kognitif siswa. Casal-Otero et al. \cite{ref2} memperkuat temuan tersebut dengan menunjukkan bahwa kurikulum literasi KB yang efektif harus memasukkan pengalaman langsung di mana siswa dihadapkan pada keluaran KB yang tidak pasti, sehingga siswa memiliki data empiris untuk mengevaluasi keluaran KB secara konkret. Dalam proyek ini, literasi KB diwujudkan melalui simulasi di mana siswa SMP berinteraksi langsung dengan keluaran klasifikasi sketsa yang bersifat probabilistik. Siswa tidak hanya melihat hasil klasifikasi, tetapi juga melihat \textit{confidence score} dan beberapa prediksi teratas (\textit{top-3 predictions}) yang memperlihatkan bahwa keluaran KB bukan jawaban tunggal yang pasti.

% ------------------------------------------------------------
\subsection{Human-in-the-Loop}
\label{subsec:hitl}
% ------------------------------------------------------------

\textit{Human-in-the-Loop} (HITL) merupakan paradigma di mana manusia dilibatkan secara aktif dalam alur kerja sistem kecerdasan buatan, khususnya pada titik-titik di mana keputusan sistem bersifat tidak pasti atau berpotensi keliru \cite{ref4}. Mosqueira-Rey et al. \cite{ref4} mendefinisikan HITL sebagai kerangka kerja yang menempatkan manusia sebagai agen evaluatif yang memiliki otoritas untuk menerima, menolak, atau mengoreksi keluaran yang dihasilkan oleh model KB. Berbeda dengan sistem KB yang sepenuhnya otonom, sistem HITL mengakui bahwa terdapat batasan fundamental pada kemampuan prediktif model dan bahwa keterlibatan manusia merupakan mekanisme esensial untuk menjaga kualitas keputusan.

Salah satu masalah utama yang menjadi rasionalisasi bagi penerapan HITL adalah \textit{automation bias}---kecenderungan manusia untuk menerima keluaran sistem otomatis tanpa evaluasi kritis yang memadai \cite{ref4}. Dalam konteks pendidikan, Memarian \& Doleck \cite{ref5} menemukan bahwa \textit{automation bias} menjadi hambatan signifikan bagi literasi KB karena mendorong siswa untuk menerima keluaran AI secara pasif. Feldman-Maggor et al. \cite{ref25} menambahkan bahwa kepercayaan guru dan siswa terhadap sistem AI EdTech sangat dipengaruhi oleh transparansi proses keputusan; tanpa mekanisme yang memungkinkan manusia menginterrogasi keluaran, keterlibatan manusia cenderung bersifat simbolis. Dalam proyek ini, mekanisme HITL diwujudkan melalui \textit{Probe UI}---antarmuka yang menghentikan sementara alur simulasi (\textit{pause-and-evaluate}) dan meminta siswa memilih antara menerima prediksi KB (\textit{Accept}), mengoreksi label (\textit{Correct}), atau menolak sepenuhnya (\textit{Override}). Keputusan siswa memiliki konsekuensi langsung dalam simulasi sehingga keterlibatan manusia bersifat substantif, bukan formalitas.

% ------------------------------------------------------------
\subsection{Explainable Interface dan Confidence Score}
\label{subsec:explainable-interface}
% ------------------------------------------------------------

\textit{Explainable Artificial Intelligence} (XAI) dalam konteks pendidikan menekankan pentingnya membuat proses dan keluaran sistem KB dapat dipahami oleh pengguna awam, termasuk siswa. Khosravi et al. \cite{ref3} mengargumentasikan bahwa transparansi sistem KB di lingkungan pendidikan bukan hanya kebutuhan teknis, melainkan prasyarat pedagogis: tanpa kemampuan untuk melihat mengapa dan seberapa yakin sistem menghasilkan keluaran tertentu, pengguna tidak memiliki dasar untuk melakukan evaluasi. Karran et al. \cite{ref8} menunjukkan bahwa visualisasi \textit{confidence score}---penyajian tingkat keyakinan model secara eksplisit---memberikan dampak signifikan terhadap cara pengguna menilai keluaran KB. Ketika \textit{confidence score} rendah, pengguna cenderung lebih berhati-hati; sebaliknya, ketika \textit{confidence score} tinggi tanpa visualisasi yang memadai, pengguna cenderung menerima keluaran tanpa pertanyaan.

Konsep kunci yang mendasari desain antarmuka dalam proyek ini adalah \textbf{ambiguitas terkontrol} (\textit{controlled ambiguity}). Sistem KB dalam simulasi sengaja dirancang agar tidak selalu menghasilkan prediksi yang sempurna, sehingga siswa memiliki kesempatan untuk mengevaluasi \cite{ref3}. Dalam praktiknya, ambiguitas terkontrol diwujudkan melalui pengaturan \textit{confidence threshold} pada setiap level simulasi: Level 1 menggunakan \textit{confidence} tinggi ($>85\%$) untuk membangun kepercayaan awal, Level 2 menurunkan \textit{confidence} ke rentang 55--70\% untuk menghadirkan ambiguitas yang membutuhkan evaluasi, dan Level 3 menempatkan \textit{confidence} di rentang 35--51\% di mana siswa hampir selalu perlu melakukan \textit{override}. Dengan demikian, desain antarmuka tidak hanya menyajikan informasi, tetapi secara sengaja mengatur derajat ketidakpastian untuk menciptakan momen-momen HITL yang bermakna.

% ------------------------------------------------------------
\subsection{Simulasi Interaktif Berlevel}
\label{subsec:simulasi-interaktif}
% ------------------------------------------------------------

Simulasi interaktif berbasis permainan (\textit{game-based learning}/GBL) telah menunjukkan potensi signifikan dalam mendukung keterlibatan aktif siswa dalam konteks pendidikan sains komputer. Videnovik et al. \cite{ref6} menemukan bahwa GBL secara konsisten meningkatkan keterlibatan siswa, terutama ketika elemen permainan dirancang untuk mendukung tujuan pembelajaran secara langsung. Lee et al. \cite{ref13} menerapkan GBL AI literacy untuk siswa \textit{upper elementary} dan menemukan bahwa pendekatan kolaboratif-inquiry berbasis permainan efektif dalam memperkenalkan konsep KB secara konkret. Gupta et al. \cite{ref14} memperluas pendekatan tersebut dengan narasi berbasis cerita (\textit{story-driven GBL}) yang terbukti mendukung siswa dalam mempelajari konsep KB yang abstrak melalui konteks naratif.

Prinsip pedagogis yang mendasari desain tiga level dalam proyek ini adalah \textbf{scaffolding}---pemberian bantuan yang melimpah pada tahap awal yang kemudian dikurangi secara bertahap seiring meningkatnya kompetensi siswa \cite{ref6}. Selaras dengan scaffolding, desain level juga didasarkan pada Teori Flow yang menyatakan bahwa individu mencapai pengalaman optimal ketika tingkat tantangan seimbang dengan tingkat keterampilan yang dimilikinya. Chan et al. \cite{ref7} secara empiris menemukan bahwa keseimbangan tantangan-keterampilan merupakan prediktor utama pengalaman \textit{flow} dalam lingkungan GBL non-kompetitif. Dalam proyek ini, Level 1 menyajikan tantangan rendah sesuai keterampilan awal siswa, Level 2 meningkatkan tantangan secara terukur seiring meningkatnya pemahaman siswa, dan Level 3 menyajikan tantangan tertinggi yang menuntut keterampilan evaluatif yang telah terbangun pada dua level sebelumnya. Kombinasi scaffolding dan Teori Flow memastikan bahwa progresi tidak hanya bertahap secara pedagogis tetapi juga terjaga secara pengalaman.

% ------------------------------------------------------------
\subsection{Pedagogical Agent --- Maskot Momo}
\label{subsec:pedagogical-agent}
% ------------------------------------------------------------

Agen pedagogis (\textit{pedagogical agent}) merupakan karakter virtual yang terintegrasi dalam lingkungan pembelajaran digital dan berfungsi memberikan panduan, umpan balik, atau motivasi kepada siswa. Schroeder, Davis \& Yang \cite{ref9} dalam tinjauan komprehensif mereka (\textit{umbrella review}) menemukan bahwa agen pedagogis dapat mendukung proses belajar jika desainnya memenuhi tiga kondisi: (1) relevansi kontekstual---agen merespons situasi spesifik yang dihadapi siswa, bukan memberikan pesan generik; (2) kesesuaian beban kognitif---agen tidak memberikan informasi yang membebani siswa secara berlebihan; serta (3) konsistensi persona---agen memiliki identitas dan gaya komunikasi yang konsisten sepanjang interaksi.

Agen pedagogis dalam proyek ini, yaitu karakter Momo, dirancang sebagai agen berbasis aturan (\textit{rule-based}) yang non-NLP, bukan sebagai \textit{chatbot} berbasis pemrosesan bahasa alami. Keputusan desain ini didasarkan pada temuan Schroeder et al. \cite{ref9} yang menunjukkan bahwa agen pedagogis yang sederhana dan kontekstual seringkali lebih efektif daripada agen yang kompleks dan \textit{overwhelming}. Hsu et al. \cite{ref15} mendukung pendekatan ini dengan menunjukkan bahwa pola perilaku siswa dalam mempelajari KB dapat dipandu secara efektif melalui umpan balik kontekstual berbasis aturan, tanpa perlu interaksi percakapan bebas. Ocak et al. \cite{ref16} memperluas temuan tersebut dengan menunjukkan bahwa analisis pola interaksi anak (\textit{embodied interactions}) dapat menjadi dasar untuk merancang agen pedagogis yang adaptif terhadap kondisi pembelajaran. Dalam implementasinya, Momo muncul pada momen-momen tertentu sesuai kondisi sistem: saat \textit{onboarding} untuk menjelaskan aturan, saat \textit{Probe UI} untuk mengarahkan perhatian siswa pada \textit{confidence score}, dan saat siswa membuat keputusan berpotensi berisiko untuk memberikan peringatan kontekstual.

% ------------------------------------------------------------
\subsection{Child-Computer Interaction untuk Siswa SMP}
\label{subsec:cci-smp}
% ------------------------------------------------------------

\textit{Child-Computer Interaction} (CCI) merupakan bidang kajian yang mempelajari bagaimana anak-anak dan remaja berinteraksi dengan teknologi digital, dengan fokus pada bagaimana karakteristik perkembangan kognitif, motorik, dan sosial memengaruhi desain interaksi yang sesuai. Dalam konteks literasi KB untuk kalangan K-12, Casal-Otero et al. \cite{ref2} menekankan bahwa strategi edukasi KB harus disesuaikan dengan tahap perkembangan usia siswa. Untuk siswa SMP yang berada pada tahap operasional formal menurut teori Piaget, siswa sudah mampu berpikir abstrak namun masih membutuhkan jembatan konkrit untuk memahami konsep yang kompleks. Wang et al. \cite{ref19} menemukan bahwa gaya kognitif siswa memengaruhi cara mereka memproses informasi dalam kursus berbasis AI, sehingga desain antarmuka harus mengakomodasi variasi tersebut. Wu \& Su \cite{ref20} mempelajari siswa kelas 5--6 dan menemukan bahwa kompleksitas antarmuka visual harus disesuaikan dengan kapasitas pemrosesan kognitif anak.

Implikasi desain dari prinsip CCI untuk siswa SMP meliputi beberapa aspek. Pertama, antarmuka harus sederhana namun tidak kekanak-kanakan---siswa SMP menolak desain yang terlalu infantil namun juga dapat mengalami \textit{choice paralysis} jika dihadapkan pada terlalu banyak pilihan \cite{ref2}. Oleh karena itu, alur simulasi dalam proyek ini dirancang secara linear tanpa percabangan, agar fokus siswa tetap pada momen keputusan HITL. Kedua, umpan balik harus bersifat non-menghukum---ketika siswa membuat keputusan yang kurang tepat, sistem memberikan kesempatan untuk mencoba lagi (\textit{retry}) alih-alih memberikan penalti, sehingga mendorong eksplorasi tanpa tekanan. Ketiga, interaksi harus memanfaatkan modalitas yang familiar bagi remaja, seperti antarmuka visual dan input berbasis kanvas, yang menjadi dasar pemilihan kanvas gambar sebagai mekanisme input utama.

% ------------------------------------------------------------
\subsection{Usability Testing}
\label{subsec:usability-testing}
% ------------------------------------------------------------

\textit{Usability testing} merupakan metode evaluasi yang mengukur sejauh mana sebuah sistem interaktif dapat digunakan secara efektif, efisien, dan memuaskan oleh pengguna target dalam konteks penggunaan tertentu. Dalam lingkup proyek ini, \textit{usability testing} dipilih sebagai metode evaluasi utama---bukan metode pre-post test yang mengklaim peningkatan kemampuan---karena tujuan evaluasi adalah mengukur keterbacaan alur interaksi, kejelasan mekanisme HITL, dan kesesuaian desain antarmuka dengan kemampuan siswa SMP. Brooke \cite{ref31} mengembangkan \textit{System Usability Scale} (SUS) sebagai instrumen \textit{quick-and-dirty} yang efektif untuk evaluasi awal sistem interaktif, dan Bangor et al. \cite{ref32} memvalidasi SUS secara empiris sebagai instrumen yang reliabel di berbagai domain aplikasi.

Aspek \textit{usability} yang diukur dalam proyek ini meliputi: (1) \textit{learnability}---seberapa mudah siswa memahami mekanisme simulasi setelah tahap \textit{onboarding}; (2) \textit{efficiency}---seberapa cepat siswa dapat menyelesaikan alur simulasi dan membuat keputusan di setiap \textit{Probe UI}; (3) \textit{error rate}---seberapa sering siswa membuat kesalahan dalam navigasi atau interpretasi \textit{confidence score}; serta (4) \textit{satisfaction}---tingkat kepuasan subjektif siswa terhadap pengalaman interaksi. Pengukuran aspek-aspek ini dilakukan melalui kombinasi observasi terstruktur, analisis log interaksi, dan kuesioner SUS pasca-simulasi. Dengan pendekatan ini, evaluasi menghasilkan data empiris mengenai kualitas desain interaksi tanpa membuat klaim yang melampaui cakupan metode yang digunakan.

% ============================================================
\section{PENELITIAN TERKAIT}
\label{sec:penelitian-terkait}
% ============================================================

Bagian ini menyajikan kajian terhadap penelitian-penelitian sebelumnya yang relevan dengan proyek \textit{Sketchbook Universe}. Setiap sub-bagian membahas temuan utama dari penelitian tersebut serta hubungannya dengan desain sistem ini.

% ------------------------------------------------------------
\subsection{AI Literacy in K-12}
\label{subsec:penelitian-ai-literacy-k12}
% ------------------------------------------------------------

Casal-Otero et al. \cite{ref2} melakukan tinjauan literatur sistematis terhadap 37 studi literasi KB pada jenjang K-12 dan menemukan bahwa strategi berbasis pengalaman langsung (\textit{experience-based}) lebih efektif dibandingkan strategi ekspositif murni. Wang \& Lester \cite{ref11} menegaskan bahwa literasi KB K-12 merupakan panggilan aksi (\textit{call to action}) yang memerlukan integrasi konsep KB ke dalam kurikulum sejak usia dini dengan pendekatan age-appropriate. Yim \& Su \cite{ref12} memperluas kajian tersebut ke konteks sekolah dasar dan menemukan bahwa alat bantu pembelajaran KB harus dirancang dengan mempertimbangkan kapasitas kognitif anak, terutama dalam memahami konsep probabilistik. Ng et al. \cite{ref1} menyediakan kerangka konseptual literasi KB yang mencakup kemampuan mengevaluasi keluaran probabilistik---aspek yang menjadi inti dari desain \textit{Sketchbook Universe}, yang menempatkan siswa dalam posisi evaluator aktif terhadap keluaran klasifikasi KB.

% ------------------------------------------------------------
\subsection{Human-in-the-Loop Machine Learning}
\label{subsec:penelitian-hitl-ml}
% ------------------------------------------------------------

Mosqueira-Rey et al. \cite{ref4} menyajikan tinjauan komprehensif terhadap \textit{state of the art} dalam \textit{Human-in-the-Loop Machine Learning} (HITL-ML) dan mengidentifikasi pola interaksi HITL yang umum, termasuk \textit{active learning}, \textit{interactive machine learning}, dan \textit{human evaluation}. Mereka secara khusus menekankan bahwa \textit{automation bias} merupakan tantangan mendasar dalam desain sistem HITL: tanpa desain antarmuka yang secara eksplisit memaksa pengguna untuk mengevaluasi keluaran, keterlibatan manusia cenderung bersifat simbolis. Memarian \& Doleck \cite{ref5} memfokuskan kajian pada penerapan HITL dalam konteks pendidikan kecerdasan buatan (AIED) dan menemukan bahwa efektivitas HITL bergantung pada kualitas informasi yang diberikan kepada manusia. Videnovik et al. \cite{ref6} mengaitkan prinsip HITL dengan pendekatan GBL dalam pendidikan sains komputer, di mana keputusan pemain memiliki konsekuensi langsung terhadap kemajuan permainan---mekanisme yang menjadi dasar desain \textit{Probe UI} pada proyek ini.

% ------------------------------------------------------------
\subsection{Game-Based Learning in Computer Science Education}
\label{subsec:penelitian-gbl-cs}
% ------------------------------------------------------------

Videnovik et al. \cite{ref6} melakukan tinjauan literatur berskala (\textit{scoping literature review}) terhadap penerapan GBL dalam pendidikan sains komputer dan mengidentifikasi bahwa GBL paling efektif ketika menerapkan \textit{consequence-driven learning}---di mana keputusan pemain memiliki konsekuensi langsung terhadap kemajuan permainan. Chan et al. \cite{ref7} mengkaji hubungan antara Teori Flow dan GBL secara empiris pada lingkungan non-kompetitif, dan menemukan bahwa keseimbangan tantangan-keterampilan merupakan prediktor yang lebih kuat terhadap pengalaman \textit{flow} dibandingkan elemen kompetisi. Lee et al. \cite{ref13} menerapkan GBL AI literacy untuk siswa \textit{upper elementary} dengan pendekatan \textit{collaborative inquiry} yang terbukti efektif dalam memperkenalkan konsep KB. Gupta et al. \cite{ref14} memperluas pendekatan tersebut dengan narasi berbasis cerita yang mendukung pemahaman konsep KB yang abstrak melalui konteks naratif. Temuan ini secara langsung mendukung desain tiga level dalam \textit{Sketchbook Universe}, yang meningkatkan tantangan evaluatif secara bertahap dari Level 1 hingga Level 3 tanpa elemen kompetisi antar pemain, serta memanfaatkan narasi visual berbasis buku sketsa sebagai konteks pembelajaran.

% ------------------------------------------------------------
\subsection{Children Behavioral Patterns in AI Learning}
\label{subsec:penelitian-behavioral-patterns}
% ------------------------------------------------------------

Catatan: Sub-bab ini menggantikan ``Real-Time Hand Gesture Monitoring'' yang DEPRECATED pasca-pivot 16/6/26 (login gesture diganti nomor absen + superadmin). Fokus kajian dialihkan ke pola perilaku anak dalam mempelajari KB, yang lebih relevan dengan desain interaksi reflektif siswa terhadap keluaran KB pada proyek ini.

Hsu et al. \cite{ref15} mengembangkan instrumen pembelajaran KB untuk anak dan mengeksplorasi pola perilaku siswa saat berinteraksi dengan sistem berbasis KB. Mereka menemukan bahwa pola perilaku siswa---seperti kecenderungan menerima keluaran sistem secara langsung versus mengevaluasi sebelum mengambil keputusan---dapat diidentifikasi melalui analisis log interaksi, dan bahwa pola ini berkorelasi dengan tingkat pemahaman siswa terhadap konsep probabilistik KB. Ocak et al. \cite{ref16} menggunakan pendekatan pengenalan pola berbasis AI untuk menganalisis interaksi \textit{embodied} anak (\textit{children's embodied interactions}) dan menemukan bahwa analisis pola spasial-temporal gerakan anak dapat memberikan indikasi tentang tingkat keterlibatan kognitif mereka. Temuan kedua penelitian ini sangat relevan dengan proyek \textit{Sketchbook Universe} karena sistem yang dikembangkan mencatat log keputusan siswa untuk analisis pola perilaku evaluatif terhadap keluaran KB. Dengan demikian, sub-bab ini memberikan landasan akademis untuk analisis data log interaksi siswa yang akan dilakukan pada tahap evaluasi.

% ------------------------------------------------------------
\subsection{State of The Art}
\label{subsec:state-of-the-art}
% ------------------------------------------------------------

Untuk memosisikan proyek \textit{Sketchbook Universe} di antara penelitian-penelitian sebelumnya, dilakukan perbandingan terhadap beberapa sistem dan kajian terkait. Tabel~\ref{tab:sota} menyajikan perbandingan berdasarkan dimensi-dimensi yang relevan dengan fokus proyek ini.

\begin{table}[H]
\centering
\caption{Perbandingan State of The Art}
\label{tab:sota}
\small
\begin{tabular}{|p{3.8cm}|c|c|c|c|c|c|}
\hline
\textbf{Sistem / Kajian} & \textbf{Literasi KB} & \textbf{HITL} & \textbf{Conf. Vis.} & \textbf{GBL Berlevel} & \textbf{Ped. Agent} & \textbf{CCI SMP} \\
\hline
Ng et al. \cite{ref1} & \checkmark & -- & -- & -- & -- & -- \\
\hline
Casal-Otero et al. \cite{ref2} & \checkmark & -- & -- & -- & -- & \checkmark \\
\hline
Mosqueira-Rey et al. \cite{ref4} & -- & \checkmark & -- & -- & -- & -- \\
\hline
Videnovik et al. \cite{ref6} & -- & -- & -- & \checkmark & -- & -- \\
\hline
Lee et al. \cite{ref13} & \checkmark & -- & -- & \checkmark & -- & -- \\
\hline
Hsu et al. \cite{ref15} & \checkmark & -- & -- & -- & \checkmark & \checkmark \\
\hline
\textbf{Sketchbook Universe} & \checkmark & \checkmark & \checkmark & \checkmark & \checkmark & \checkmark \\
\hline
\end{tabular}
\end{table}

Berdasarkan Tabel~\ref{tab:sota}, terlihat bahwa tidak ada satupun sistem atau kajian sebelumnya yang mengintegrasikan keenam dimensi secara bersamaan. Ng et al. \cite{ref1} dan Casal-Otero et al. \cite{ref2} membahas literasi KB secara konseptual tanpa mengimplementasikan mekanisme HITL atau visualisasi \textit{confidence}. Mosqueira-Rey et al. \cite{ref4} mengkaji HITL secara komprehensif namun tidak membahas implementasi untuk kalangan K-12. Videnovik et al. \cite{ref6} dan Lee et al. \cite{ref13} meneliti GBL dan progresi berlevel namun tidak mengintegrasikannya dengan mekanisme HITL untuk evaluasi keluaran KB. Hsu et al. \cite{ref15} membahas pola perilaku anak belajar KB tanpa implementasi sistem interaktif yang konkret.

Proyek \textit{Sketchbook Universe} secara unik mengintegrasikan keenam dimensi tersebut: (1) literasi KB melalui pengalaman langsung berinteraksi dengan keluaran klasifikasi probabilistik; (2) mekanisme HITL melalui \textit{Probe UI} yang menghentikan simulasi untuk evaluasi; (3) visualisasi \textit{confidence score} yang memperlihatkan tingkat keyakinan model; (4) progresi berlevel berdasarkan prinsip scaffolding dan Teori Flow; (5) agen pedagogis Momo yang memberikan umpan balik kontekstual; serta (6) desain interaksi CCI yang sesuai dengan kemampuan kognitif siswa SMP. Integrasi ini merupakan kontribusi utama proyek terhadap ruang kajian yang ada, memposisikan \textit{Sketchbook Universe} sebagai sistem yang menjembatani kesenjangan antara teori literasi KB, mekanisme HITL, dan desain interaksi untuk siswa SMP.
%% END BAB 2 CONTENT %%

% ============================================================
% DAFTAR PUSTAKA (subset — hanya ref yang disitir di Bab 2)
% Daftar pustaka lengkap (32 refs) ada di daftar_pustaka_can_v1.tex
% ============================================================
\chapter*{DAFTAR PUSTAKA}
\addcontentsline{toc}{chapter}{DAFTAR PUSTAKA}

\begin{thebibliography}{99}

\bibitem{ref1}
Ng, D.T.K., Leung, J.K.L., Chu, S.K.W., \& Qiao, M.S. (2021). Conceptualizing AI literacy: An exploratory review. \textit{Computers and Education: Artificial Intelligence}, 2, 100041.

\bibitem{ref2}
Casal-Otero, V., Catala, A., Fern\'andez-Morante, C., Taboada, M., Caeiro-Rodr\'iguez, M., \& Barro, S. (2023). AI literacy in K-12: A systematic literature review. \textit{International Journal of STEM Education}, 10, 29.

\bibitem{ref3}
Khosravi, H., Shum, S.B., Chen, G., Conati, C., Tsai, Y.-S., Kay, J., Knight, S., Martinez-Maldonado, R., Sadat-Milani, L., \& Ga\v{s}evi\'c, D. (2022). Explainable Artificial Intelligence in education. \textit{Computers and Education: Artificial Intelligence}, 3, 100074.

\bibitem{ref4}
Mosqueira-Rey, E., Hern\'andez-Pereira, E., Alonso-R\'ios, D., Bobillo-Ares, B., \& Moret-Bonillo, V. (2023). Human-in-the-loop machine learning: A state of the art. \textit{Artificial Intelligence Review}, 56(4), 3005--3054.

\bibitem{ref5}
Memarian, B., \& Doleck, T. (2024). Human-in-the-loop in artificial intelligence in education: A review and entity-relationship (ER) analysis. \textit{Computers in Human Behavior: Artificial Humans}, 2, 100053.

\bibitem{ref6}
Videnovik, M., Vlahovi\v{c}-Steti\'c, V., Kišiček, S., \& Vošner, H.B. (2023). Game-based learning in computer science education: A scoping literature review. \textit{International Journal of STEM Education}, 10, 54.

\bibitem{ref7}
Chan, K.K., Wan, K.M., \& King, R.B. (2021). Performance over enjoyment? Effect of game-based learning on learning outcome and flow experience. \textit{Frontiers in Education}, 6, 660376.

\bibitem{ref8}
Karran, A.J., Demazure, T., \& Hudelot, C. (2022). Designing for confidence: The impact of visualizing artificial intelligence decisions. \textit{Frontiers in Neuroscience}, 16.

\bibitem{ref9}
Schroeder, N.L., Davis, C.S., \& Yang, Y. (2025). Designing and learning with pedagogical agents: An umbrella review. \textit{Journal of Educational Computing Research}, 62(8).

\bibitem{ref11}
Wang, N., \& Lester, J. (2023). K-12 education in the age of AI: A call to action for K-12 AI literacy. \textit{International Journal of Artificial Intelligence in Education}, 33(3), 631--640.

\bibitem{ref12}
Yim, I. H. Y., \& Su, J. (2025). Artificial intelligence literacy education in primary schools: A review. \textit{International Journal of Technology and Design Education}, 35, Article 57.

\bibitem{ref13}
Lee, S., Mott, B., Ottenbreit-Leftwich, A., Scribner, A., Taylor, S., Glazewski, K., \& Lester, J. (2021). AI-infused collaborative inquiry in upper elementary school: A game-based learning approach. \textit{Proceedings of the AAAI Conference on Artificial Intelligence}, 35(17), 15653--15661.

\bibitem{ref14}
Gupta, A., Lee, S., Mott, B., Chakraburty, S., Scribner, A., Glazewski, K., Ottenbreit-Leftwich, A., \& Lester, J. (2024). Supporting upper elementary students in learning AI concepts with story-driven game-based learning. \textit{Proceedings of the AAAI Conference on Artificial Intelligence}, 38(21), 23321--23329.

\bibitem{ref15}
Hsu, T. C., Abelson, H., Lao, N., Tseng, Y. H., \& Lin, Y. T. (2021). Behavioral-pattern exploration and development of an instructional tool for young children to learn AI. \textit{Computers and Education: Artificial Intelligence}, 2, Article 100060.

\bibitem{ref16}
Ocak, C., Kopcha, T. J., \& Dey, R. (2023). An AI-enhanced pattern recognition approach to temporal and spatial analysis of children's embodied interactions. \textit{Computers and Education: Artificial Intelligence}, 4, Article 100155.

\bibitem{ref19}
Wang, C. J., Zhong, H. X., Chiu, P. S., Chang, J. H., \& Chen, M. Y. (2022). Research on the impacts of cognitive style and computational thinking on college students in a visual artificial intelligence course. \textit{Frontiers in Psychology}, 13, Article 864416.

\bibitem{ref20}
Wu, S. Y., \& Su, Y. S. (2021). Visual programming environments and computational thinking performance of fifth-and sixth-grade students. \textit{Journal of Educational Computing Research}, 59(7), 1413--1441.

\bibitem{ref25}
Feldman-Maggor, Y., Cukurova, M., Kent, C., Luckin, R., \& Baur-Wolfson, M. (2025). The impact of explainable AI on teachers' trust and acceptance of AI EdTech recommendations. \textit{International Journal of Artificial Intelligence in Education}, 35(1), 1--25.

\bibitem{ref31}
Brooke, J. (1996). SUS: A `Quick and Dirty' Usability Scale. In \textit{Usability Evaluation in Industry} (pp. 189--194). Taylor \& Francis.

\bibitem{ref32}
Bangor, A., Kortum, P.T., \& Miller, J.T. (2008). An empirical evaluation of the System Usability Scale. \textit{International Journal of Human-Computer Interaction}, 24(6), 574--594.

\end{thebibliography}

\end{document}
